% ---- pg-363 （书页 351） ----

化主要由市场活动和与之不相关的行业活动造成. 股票价格行为的这种解释也是 King 所主张的(见文献[12]).

其他主成分不易解释, 总体上表示各个股票的特性. 无论如何, 它们对解释样本的总方差没有什么帮助. \hfill$\blacksquare$

\noindent\textbf{例 8.6}\quad (来自有特殊结构的相关矩阵的成分)

遗传学常常涉及特征遗传, 对这些特征可在动物的生命周期内做数次测量. 对 $n=150$ 只雌鼠, 在它们生下前 4 只幼鼠后都立即称出雌鼠的体重(单位克).\footnote{{\footnotesize 数据经 J. J. 拉特利奇(Rutledge)许可使用.}}样本均值向量和样本相关矩阵分别为
\[
\bar{\mathbf{x}}'=[39.88,45.08,48.11,49.95]
\]
和
\[
\mathbf{R}=\begin{bmatrix}
1.000&0.750\ 1&0.632\ 9&0.636\ 3\\
0.750\ 1&1.000&0.692\ 5&0.738\ 6\\
0.632\ 9&0.692\ 5&1.000&0.662\ 5\\
0.636\ 3&0.738\ 6&0.662\ 5&1.000
\end{bmatrix}
\]
该矩阵的特征值是
\[
\hat{\lambda}_1=3.085,\qquad \hat{\lambda}_2=0.382,\qquad \hat{\lambda}_3=0.342\qquad\text{和}\qquad \hat{\lambda}_4=0.217
\]
我们注意到第一特征值近似等于 $1+(p-1)\bar{r}=1+(4-1)\times 0.6854=3.056$, 其中 $\bar{r}$ 是 $\mathbf{R}$ 中非对角元的算术平均值. 余下的特征值都很小并且差不多相等, 虽然 $\hat{\lambda}_4$ 比 $\hat{\lambda}_2$ 和 $\hat{\lambda}_3$ 更小一点. 这样就有证据说, 对应的总体相关矩阵 $\boldsymbol{\rho}$ 可以是\hyperref[eq:8-15]{式(8-15)}的“等相关”形式. 这一概念在例 8.9 将进一步探讨.

第一主成分
\[
\hat{y}_1=\hat{\mathbf{e}}_1'\mathbf{z}=0.49z_1+0.52z_2+0.49z_3+0.50z_4
\]
所解释的总方差占 $100(\hat{\lambda}_1/p)\%=100\times(3.058/4)\%=76\%$. 虽然平均产后体重随时间增加, 有(近似)相等系数的第一主成分仍然相当清楚地解释了体重上的变差. \hfill$\blacksquare$

\mvsub{评注}\quad 从样本方差或是从相关矩阵得到的最后一个特征值, 若特别小, 它可能显露出数据集合中一个未被注意到的相依性. 若真的出现这种情况, 一个或多个变量是多余的, 应该除去. 考虑这一个情况: $x_1, x_2$ 和 $x_3$ 是测验得分, 总成绩 $x_4=x_1+x_2+x_3$. 此时, 虽然线性组合 $\mathbf{e}_4'\mathbf{x}=[1,1,1,-1]\mathbf{x}=x_1+x_2+x_3-x_4$ 总是零, 但计算特征值时一个舍入误差可能导致一个非零的微小数值. 如果 $x_4$ 与 $(x_1,x_2,x_3)$ 的线性表达式一开始时被忽略了, 那么最小特征值-特征向量对, 将提供存在这个线性关系的一个线索(见 3.4 节中的讨论).

这样一来, 虽然“大”特征值和它相应的特征向量, 在主成分分析时是重要的, 但很接近于 0 的特征值通常也不应该忽略. 与后面这些的特征值相联系的特征向量, 可能指出数据集合中存在的线性相依关系, 而这在随后的分析中可能引出解释和计算方面的问题.

\section{主成分的图形表示}  % 原书 8.4
\phantomsection\label{sec:8.4}

主成分的图形表示除了能核查正态性假定, 还能揭示可疑的观测值. 由于主成分是原变量的线性组合, 因此期望它们是近似正态就不是没有道理的. 当要把前几个主成分作为输入
